%% toc: P2 17 17.8
\section{Proof of the quantified subformula property}
\label{sec:p02-w02b-p02-020-subformula-proof}

\emph{Purpose.}  \Cref{sec:p02-fol-subformula-property} defines the
first-order subformula-instance relation
(\cref{def:p02-fol-subformula-instance}) and states that every formula in a
derivation on \texttt{sheet:lk} without Cut or Mix is
such an instance of a formula in the endsequent
(\cref{thm:p02-fol-subformula-property}).  This section proves that claim in a
stronger, traced form: each formula occurrence is followed downward through
the derivation, one rule at a time, to a definite occurrence in the endsequent,
and each step is recorded.  The proof also shows exactly where the first-order
property departs from literal subformulahood and why Cut, Mix and the
equality rule of \texttt{sheet:lk-eq} lie outside it.

\emph{Inputs.}  A first-order language \(\Lang\) with terms
\(\operatorname{Term}(\Lang)\) (\cref{sec:fol-syntax-semantics-01}) and formulas
\(\operatorname{Fm}(\Lang)\) (\cref{sec:fol-syntax-semantics-02}); the total
capture-avoiding substitution \(\varphi[t/x]\) of
\cref{def:p02-fol-capture-avoiding-substitution}, which is the operation used
by the quantifier rules of the sheet; and a finite derivation tree on
\texttt{sheet:lk}: two-sided sequents of finite
multisets, contexted initial sequents \(\Gamma,\varphi\Rightarrow\varphi,\Delta\)
and \(\Gamma,\bot\Rightarrow\Delta\), weakening, contraction, the rules for
\(\neg,\land,\lor,\to\), and the four quantifier rules, with no equality rule.
Formulas are compared literally, not modulo alpha-equivalence.

\emph{Conclusions.}  Every formula occurrence of a derivation containing
neither Cut nor Mix reaches a unique endsequent occurrence by a finite chain of
steps, each of which keeps the formula, passes to an immediate propositional
subformula, or passes to a term instance of a quantified formula
(\cref{thm:p02-w02b-p02-020-traced}).  Hence it is a first-order subformula
instance of that endsequent formula; this is the assertion of
\cref{thm:p02-fol-subformula-property}.  When no endsequent formula contains a
quantifier, every formula of the derivation is a literal subformula of an
endsequent formula (\cref{thm:p02-w02b-p02-020-quantifier-free}).

\subsection{Instance sets, literal subformulas and successors}
\label{sec:p02-w02b-p02-020-conventions}

\begin{hbconvention}{def:p02-w02b-p02-020-instance-set}{Instance set and literal subformulas}
The \emph{immediate propositional subformula} of \(\neg\varphi\) is
\(\varphi\); those of \(\varphi\circ\psi\), for
\(\circ\in\{\land,\lor,\to\}\), are \(\varphi\) and \(\psi\).  Atomic formulas,
\(\bot\) and quantified formulas have none.  These are well defined by unique
readability (\cref{thm:fol-unique-readability}).  A formula is
\emph{quantifier-free} when neither \(\forall\) nor \(\exists\) occurs in it.

For a formula \(\chi\), the \emph{instance set} \(\operatorname{SI}(\chi)\) is
the least set of formulas that contains \(\chi\) and is closed under
\begin{enumerate}
  \item passing from a formula to each of its immediate propositional
        subformulas, and
  \item passing from \(Qx\,\varphi\), for \(Q\in\{\forall,\exists\}\), to
        \(\varphi[t/x]\) for every \(t\in\operatorname{Term}(\Lang)\).
\end{enumerate}
The terms in clause~2 are those of \(\Lang\), variables included; a term of
\(\Lang\) is what the clause ``every term \(t\) in the language of the
derivation'' of \cref{def:p02-fol-subformula-instance} ranges over.  The
\emph{literal subformulas} of \(\chi\) form the least set that contains
\(\chi\) and is closed under clause~1 and under passing from \(Qx\,\varphi\)
to \(\varphi\) itself, with \(x\) left free.
\index{subformula!first-order instance}
\index{subformula!literal}
\end{hbconvention}

The set \(\operatorname{SI}(\chi)\) is the section at \(\chi\) of the relation
of \cref{def:p02-fol-subformula-instance}: the formulas \(\psi\) that are
first-order subformula instances of \(\chi\) are exactly the members of
\(\operatorname{SI}(\chi)\) (\cref{thm:p02-w02b-p02-020-closure}).  It can be infinite: for a unary
relation symbol \(P\), \(\operatorname{SI}(\forall x_0\,P(x_0))\) contains
\(P(x_0)[t/x_0]=P(t)\) for every term \(t\), and distinct terms give distinct
formulas.

\begin{hbconvention}{def:p02-w02b-p02-020-successor}{Occurrences and their successors}
A rule instance on \texttt{sheet:lk} is given together
with the decomposition of each premise and of the conclusion into the context
multisets and displayed formulas written in the rule.  When a formula occurs
several times in one sequent, the occurrences are kept apart by this
decomposition and by a fixed one-to-one matching between each context multiset
of a premise and the same multiset in the conclusion.  For every rule other
than Cut and Mix, each formula occurrence in a premise has exactly one
\emph{successor} in the conclusion:
\begin{enumerate}
  \item an occurrence in a context multiset \(\Gamma\) or \(\Delta\) of a
        premise has as successor the matched occurrence of the same formula in
        the same multiset of the conclusion (both premises of
        \(\land R\), \(\lor L\) and \(\to L\) carry the context
        \(\Gamma,\Delta\), and each of the two matched occurrences has the
        conclusion's occurrence as successor);
  \item in \(C_L\) and \(C_R\), each of the two displayed occurrences of
        \(\varphi\) has as successor the displayed \(\varphi\) of the
        conclusion;
  \item in a logical rule (\(\neg L\), \(\neg R\), \(\land L\), \(\land R\),
        \(\lor L\), \(\lor R\), \(\to L\), \(\to R\), \(\forall L\),
        \(\forall R\), \(\exists L\), \(\exists R\)), each displayed premise
        formula has as successor the displayed compound formula of the
        conclusion, called the \emph{principal} formula.
\end{enumerate}
Weakening premises consist of context occurrences only, and the initial
sequents \(\mathrm{id}\) and \(\bot L\) have no premises.  Endsequent
occurrences have no successor.  Following successors from an occurrence gives
its \emph{successor chain}.  This is the immediate-descendant relation of
\cite[\S1.2.3, p.~12]{Buss1998Intro}, transferred to the rules of this sheet.
\index{successor of a formula occurrence}
\index{sequent calculus!formula occurrence}
\end{hbconvention}

\subsection{Results}
\label{sec:p02-w02b-p02-020-results}

\begin{hblemma}{thm:p02-w02b-p02-020-closure}{proved=annex:p02-w02b-p02-020-closure}{Instance sets, chains and transitivity}
For all formulas \(\psi,\varphi,\chi\):
\begin{enumerate}
  \item \(\psi\) is a first-order subformula instance of \(\chi\) in the sense
        of \cref{def:p02-fol-subformula-instance} if and only if
        \(\psi\in\operatorname{SI}(\chi)\);
  \item \(\psi\in\operatorname{SI}(\chi)\) if and only if there is a finite
        sequence \(\chi=\psi_0,\psi_1,\ldots,\psi_n=\psi\) with \(n\ge0\) in
        which each \(\psi_{i+1}\) is an immediate propositional subformula of
        \(\psi_i\), or \(\psi_i=Qx\,\theta\) and
        \(\psi_{i+1}=\theta[t/x]\) for some \(Q\in\{\forall,\exists\}\) and
        \(t\in\operatorname{Term}(\Lang)\);
  \item if \(\varphi\in\operatorname{SI}(\chi)\), then
        \(\operatorname{SI}(\varphi)\subseteq\operatorname{SI}(\chi)\); in
        particular \(\psi\in\operatorname{SI}(\varphi)\) and
        \(\varphi\in\operatorname{SI}(\chi)\) give
        \(\psi\in\operatorname{SI}(\chi)\).
\end{enumerate}
\end{hblemma}

\begin{hblemma}{thm:p02-w02b-p02-020-literal}{proved=annex:p02-w02b-p02-020-literal}{Literal subformulas are instances}
\begin{enumerate}
  \item For every formula \(\varphi\) and variable \(x\),
        \(\varphi[x/x]=\varphi\).
  \item Every literal subformula of \(\chi\) belongs to
        \(\operatorname{SI}(\chi)\).
  \item If \(\chi\) is quantifier-free, then \(\operatorname{SI}(\chi)\) is
        the set of literal subformulas of \(\chi\), and every member of it is
        quantifier-free.
\end{enumerate}
\end{hblemma}

The inclusion in clause~2 can be proper once \(\chi\) contains a quantifier:
for a unary relation symbol \(P\),
\(\operatorname{SI}(\forall x_0\,P(x_0))\) contains \(P(x_1)\), which is not a
literal subformula of \(\forall x_0\,P(x_0)\), whose literal subformulas are
\(\forall x_0\,P(x_0)\) and \(P(x_0)\).

\begin{hbtheorem}{thm:p02-w02b-p02-020-traced}{proved=annex:p02-w02b-p02-020-traced}{Traced subformula-instance property on \(\mathsf{LK}\) without Cut or Mix}
Let \(\mathcal D\) be a derivation on
\texttt{sheet:lk} in which neither Cut nor Mix occurs.
For every formula occurrence \(o\) of \(\mathcal D\):
\begin{enumerate}
  \item the successor chain \(o=o_0,o_1,\ldots,o_k\) of \(o\) is finite and
        ends at an occurrence \(o_k\) of the endsequent of \(\mathcal D\);
  \item for each \(i<k\), writing \(\psi_i\) for the formula at \(o_i\),
        either \(\psi_i=\psi_{i+1}\), or \(\psi_i\) is an immediate
        propositional subformula of \(\psi_{i+1}\), or
        \(\psi_{i+1}=Qx\,\theta\) and \(\psi_i=\theta[t/x]\) with
        \(t\in\operatorname{Term}(\Lang)\); in the last case \(t\) is the term
        of the rule instance when the step is made by \(\forall L\) or
        \(\exists R\), and the eigenvariable \(y\) when it is made by
        \(\forall R\) or \(\exists L\);
  \item the formula at \(o\) belongs to \(\operatorname{SI}(\psi_k)\).
\end{enumerate}
In particular, every formula occurrence of \(\mathcal D\) is a first-order
subformula instance of some formula occurrence in the endsequent, which is the
assertion of \cref{thm:p02-fol-subformula-property}, including its term
instances from \(\forall L\) and \(\exists R\) and its eigenvariable instances
from \(\forall R\) and \(\exists L\).
\end{hbtheorem}

\begin{hbcorollary}{thm:p02-w02b-p02-020-quantifier-free}{proved=annex:p02-w02b-p02-020-quantifier-free}{Quantifier-free endsequents give literal subformulas}
If \(\mathcal D\) is a derivation on \texttt{sheet:lk}
containing neither Cut nor Mix, and no formula of its endsequent contains a
quantifier, then every formula occurring in \(\mathcal D\) is quantifier-free
and is a literal subformula of a formula of the endsequent.  In particular no
quantifier rule occurs in \(\mathcal D\).
\end{hbcorollary}

\emph{Assumptions used.}  The proof of
\cref{thm:p02-w02b-p02-020-traced} uses three facts about the derivation: no
Cut or Mix occurs; every rule is one of the listed rules of
\texttt{sheet:lk}, each with a premise shape in which every
premise formula is a context formula, a contracted formula, an immediate
propositional subformula of the principal formula, or an instance
\(\theta[t/x]\) of the principal formula \(Qx\,\theta\); and the substitution in
the quantifier rules is the same operation as in clause~2 of
\cref{def:p02-w02b-p02-020-instance-set}.  The eigenvariable conditions of
\(\forall R\) and \(\exists L\) are part of the definition of a derivation on
the sheet, and the theorem is stated for such derivations.  The argument does
not appeal to them: an instance \(\theta[y/x]\) lies in
\(\operatorname{SI}(Qx\,\theta)\) whether or not \(y\) is fresh.  Freshness
matters for validity instead.  Without it, \(\forall R\) with
\(\varphi=P(x_1)\), \(x=y=x_1\) and \(\Gamma=\{P(x_1)\}\) would take the
initial sequent \(P(x_1)\Rightarrow P(x_1)\) to
\(P(x_1)\Rightarrow\forall x_1\,P(x_1)\), which fails in a two-element
structure where \(P\) holds of exactly the element assigned to \(x_1\); every
formula of that two-line tree is still in
\(\operatorname{SI}(\forall x_1\,P(x_1))\).

\subsection{Worked example}
\label{sec:p02-w02b-p02-020-example}

\begin{hbexample}{def:p02-w02b-p02-020-example}{An instance with a renamed bound variable}
Let \(R\) be a binary relation symbol, and use the variable order
\(x_0,x_1,x_2,\ldots\) of
\cref{def:p02-fol-capture-avoiding-substitution}.  The following is a
derivation on \texttt{sheet:lk}, written from the top,
each line obtained from the line above it.
\begin{enumerate}
  \item \(R(x_1,x_1)\Rightarrow R(x_1,x_1)\), by \(\mathrm{id}\) with
        \(\Gamma\) and \(\Delta\) empty.
  \item \(\forall x_2\,R(x_1,x_2)\Rightarrow R(x_1,x_1)\), by \(\forall L\)
        with \(\varphi=R(x_1,x_2)\), \(x=x_2\) and \(t=x_1\), since
        \(R(x_1,x_2)[x_1/x_2]=R(x_1,x_1)\).
  \item \(\forall x_0\forall x_1\,R(x_0,x_1)\Rightarrow R(x_1,x_1)\), by
        \(\forall L\) with \(\varphi=\forall x_1\,R(x_0,x_1)\), \(x=x_0\) and
        \(t=x_1\).  The substitution renames the binder: the bound variable
        \(x_1\) differs from \(x_0\) and occurs in \(t\), so the third clause
        of \cref{def:p02-fol-capture-avoiding-substitution} applies, and the
        least variable outside \(\{x_0,x_1\}\) is \(x_2\).  Thus
        \((\forall x_1\,R(x_0,x_1))[x_1/x_0]=\forall x_2\,R(x_1,x_2)\), the
        premise formula of this step.
  \item \(\forall x_0\forall x_1\,R(x_0,x_1)\Rightarrow\forall x_1\,R(x_1,x_1)\),
        by \(\forall R\) with \(\varphi=R(x_1,x_1)\), \(x=x_1\) and
        eigenvariable \(y=x_1\), since \(R(x_1,x_1)[x_1/x_1]=R(x_1,x_1)\).
        Both formulas of the conclusion are sentences, so \(x_1\) occurs free
        in no formula of \(\Gamma,\Delta,\forall x_1\,R(x_1,x_1)\).
\end{enumerate}
The formula occurrences, their successor chains and the witnesses required by
\cref{thm:p02-w02b-p02-020-traced} are as follows.
\begin{itemize}
  \item \(\forall x_2\,R(x_1,x_2)\) on line~2 has successor
        \(\forall x_0\forall x_1\,R(x_0,x_1)\) on line~3, of which it is the
        instance at \(t=x_1\).  It is not a literal subformula of either
        endsequent formula: the literal subformulas of the left formula are
        \(\forall x_0\forall x_1\,R(x_0,x_1)\), \(\forall x_1\,R(x_0,x_1)\)
        and \(R(x_0,x_1)\), and those of the right formula are
        \(\forall x_1\,R(x_1,x_1)\) and \(R(x_1,x_1)\).  Its bound variable
        \(x_2\) occurs nowhere in the endsequent.
  \item The antecedent \(R(x_1,x_1)\) on line~1 has successors
        \(\forall x_2\,R(x_1,x_2)\) and then
        \(\forall x_0\forall x_1\,R(x_0,x_1)\); its chain is two instance steps.
  \item The succedent \(R(x_1,x_1)\) on lines 1 to 3 is a context occurrence
        until line~4, where it is the eigenvariable instance of the principal
        formula \(\forall x_1\,R(x_1,x_1)\).
\end{itemize}
The naive replacement of \(x_0\) by \(x_1\) in \(\forall x_1\,R(x_0,x_1)\)
would give \(\forall x_1\,R(x_1,x_1)\), in which the inserted \(x_1\) is
captured.  The relation of \cref{def:p02-fol-subformula-instance} uses the
capture-avoiding operation, so it records the renamed formula
\(\forall x_2\,R(x_1,x_2)\) that the rule actually introduces.
\end{hbexample}

\subsection{Failure cases and limitations}
\label{sec:p02-w02b-p02-020-limits}

\begin{hbwarning}{def:p02-w02b-p02-020-cut-failure}{A Cut formula need not be an instance}
Let \(P\) and \(Q\) be distinct zero-ary relation symbols.  The sequents
\(P\Rightarrow P,Q\) and \(Q,P\Rightarrow P\) are instances of
\(\mathrm{id}\).  Cut on \(Q\) gives \(P,P\Rightarrow P,P\), and \(C_L\)
followed by \(C_R\) gives \(P\Rightarrow P\).  The formula \(Q\) occurs in this
derivation but \(\operatorname{SI}(P)=\{P\}\), since an atomic formula has no
immediate propositional subformula and is not quantified.  So the conclusion of
\cref{thm:p02-w02b-p02-020-traced} fails once Cut is admitted; the Cut
formulas have no successor in \cref{def:p02-w02b-p02-020-successor}, and that
is the point at which the proof breaks.  The same derivation with Mix in place
of Cut, with \(m=n=1\), fails in the same way.
\end{hbwarning}

\begin{hbwarning}{def:p02-w02b-p02-020-equality-failure}{The equality rule is outside the theorem}
In the equality calculus \texttt{sheet:lk-eq}, the replacement rule \(=L\)
passes from the premise \(\Gamma,\varphi[r/x]\Rightarrow\Delta\) to
\(\Gamma,r=t,\varphi[t/x]\Rightarrow\Delta\).  Take distinct constants \(a,b\),
a unary relation symbol \(P\) and a zero-ary relation symbol \(Q\).  From the
initial sequent \(P(a),Q\Rightarrow Q\), the rule with \(\varphi=P(x_0)\),
\(x=x_0\), \(r=a\) and \(t=b\) gives \(Q,a=b,P(b)\Rightarrow Q\).  The two-step derivation
contains neither Cut nor Mix.  Its formula \(P(a)\) is not in
\(\operatorname{SI}(Q)\), \(\operatorname{SI}(a=b)\) or
\(\operatorname{SI}(P(b))\), each of which is the singleton of its atomic
formula.  The equality-free hypothesis of
\cref{thm:p02-fol-subformula-property} is therefore needed as stated.
\end{hbwarning}

\emph{Limitations.}  The theorem is a statement about derivations that already
contain neither Cut nor Mix.  It does not say which sequents have such
derivations on \texttt{sheet:lk}; that is a cut-elimination
or cut-free completeness question, and the registry note on the sheet records
that the Part-7 cut-elimination sketch is not a theorem about this
presentation.  The bound on formulas is not a finiteness bound: clause~2
ranges over all terms, so an instance set can be infinite, as
\(\operatorname{SI}(\forall x_0\,P(x_0))\) is, and the property restricts the
shape of the formulas in a derivation but not the terms that instantiate
quantifiers.  The
contexted initial sequents may contain arbitrary formulas, including quantified
ones, in \(\Gamma\) and \(\Delta\); such formulas are carried down as context
and are covered by clause~1 of the successor convention.

\subsection{Relationships and sources}
\label{sec:p02-w02b-p02-020-relations}

For the defined-negation calculus \texttt{sheet:lk-defneg}, cut elimination
is \cref{thm:p07-classical-cut-elimination}, and
\cref{thm:p07-classical-analytic} records the consequences: literal
subformulas, after expanding defined negation, for propositional sequents, and
for first-order cut-free proofs the closure of the end formulas under immediate
subformulas, capture-avoiding term instances and alpha-renaming.  Those results
concern a different rule inventory and do not transfer to the primitive-negation
\texttt{sheet:lk}.  On \texttt{sheet:lk} the literal form is recovered by
\cref{thm:p02-w02b-p02-020-quantifier-free} whenever the endsequent is
quantifier-free.  For the intuitionistic sheet the analogous first-order
statement, taken modulo alpha-equivalence, is \cref{thm:p07-lj-analyticity}
with \cref{def:p07-lj-instance-subformula}.  The present theorem needs no
alpha-renaming clause and compares formulas literally, because the quantifier
rules of \texttt{sheet:lk} use exactly the deterministic substitution that
appears in the instance set, and every other rule passes formulas down
unchanged or to their principal formula.

Buss gives the ancestor argument for the cut-free propositional system
\(\mathsf{PK}\) and states the resulting subformula property
\cite[\S1.2.3, p.~12, and Prop.~1.2.4, p.~13]{Buss1998Intro}.  His first-order
calculus \(\mathsf{LK}\) separates free from bound variables, so that a
substituted term is always free for the variable it replaces; under that
convention a subformula of a formula is in general a semiformula, and he
remarks that without the convention the cut-elimination theorem would need a
slight reformulation \cite[\S2.3.1, pp.~31--32]{Buss1998Intro}.  His
quantifier rules take an arbitrary term in \(\forall{:}\mathrm{left}\) and
\(\exists{:}\mathrm{right}\) and an eigenvariable that must not appear in the
lower sequent's side formulas in \(\forall{:}\mathrm{right}\) and
\(\exists{:}\mathrm{left}\) \cite[Def.~2.3.2, p.~32]{Buss1998Intro}.  The
present sheet has one class of variables and total capture-avoiding
substitution instead, which is why the instance set is closed under
\(\theta[t/x]\) for every term and may contain renamed binders, as in
\cref{def:p02-w02b-p02-020-example}.  The proof written here is for the rules
of \texttt{sheet:lk} as registered.
